\chapter*{Résumé}
\addcontentsline{toc}{chapter}{Résumé}

Dans ce travail de fin d'études, on s'est penché sur la prédiction des chances de réussite aux concours d'entrée à l'université. On est parti d'un jeu de données avec différentes variables académiques - les scores GRE et TOEFL, la moyenne générale CGPA, et quelques autres critères qui peuvent jouer.

Le boulot a consisté d'abord à nettoyer et préparer ces données, puis à tester plusieurs approches de modélisation. Finalement, on a opté pour une forêt aléatoire (Random Forest) qui semblait donner les meilleurs résultats pour prédire l'admission des candidats.

Les résultats sont plutôt bons, avec une précision qui dépasse 90% sur nos tests. Ça montre que les critères qu'on a choisis sont effectivement pertinents pour ce type de prédiction. À l'avenir, ce travail pourrait servir de base pour développer des outils d'aide à la décision, autant pour les étudiants que pour les établissements.

\chapter*{Abstract}
\addcontentsline{toc}{chapter}{Abstract}

This final year project tackles the prediction of success rates in university entrance exams. We worked with a dataset containing various academic indicators including GRE and TOEFL scores, CGPA, and other relevant factors that might influence admission chances.

The work involved cleaning and preparing the data, then experimenting with different modeling approaches. We ultimately selected a Random Forest classifier which showed the best performance in predicting candidate admission.

The results are quite encouraging, with accuracy exceeding 90% in our tests. This demonstrates that the chosen variables are indeed relevant for this type of prediction. In the future, this work could serve as a foundation for developing decision support tools for both students and educational institutions.